\documentclass[11pt,a4paper,openany]{report}
% preamble.tex -- packages and macros of the SSCDO#2 porting report.
% Macros are defined here only; check.py parses \misura, \misurau, \studio,
% \scheda, \prova, \questione by name and arity.
\usepackage{fontspec}
\usepackage{unicode-math}
\setmainfont{Libertinus Serif}
\setsansfont{Libertinus Sans}
\setmonofont[Scale=MatchLowercase]{Libertinus Mono}
\setmathfont{Libertinus Math}
\usepackage[italian]{babel}
\usepackage[a4paper,margin=2.3cm,marginparwidth=1.8cm,marginparsep=3mm]{geometry}
\usepackage{amsmath,booktabs,longtable,array,tabularx,enumitem}
\usepackage{siunitx}
\sisetup{output-decimal-marker={,},group-separator={\,},group-minimum-digits=5,
  retain-explicit-plus=true,range-phrase={--},per-mode=symbol}
\usepackage{xcolor}
\definecolor{sdeciso}{HTML}{2E7D32}
\definecolor{sdaprovare}{HTML}{EF6C00}
\definecolor{sdomanda}{HTML}{1565C0}
\usepackage{tikz}
\usetikzlibrary{positioning,arrows.meta}
\usepackage{listings}
\lstdefinestyle{faust}{basicstyle=\ttfamily\small,breaklines=true,
  keepspaces=true,columns=fullflexible,frame=single,framesep=4pt}
\lstset{style=faust}
\usepackage{tcolorbox}
\tcbuselibrary{breakable}
\usepackage[hidelinks]{hyperref}

% ---- state tags
\newcommand{\statotag}[2]{\textcolor{#1}{\textsf{\small\bfseries #2}}}
\newcommand{\deciso}{\statotag{sdeciso}{DECISO}}
\newcommand{\daprovare}{\statotag{sdaprovare}{DA PROVARE}}
\newcommand{\domanda}{\statotag{sdomanda}{DOMANDA}}

% ---- numbers and sources
\newcommand{\misura}[2]{\num{#1}}
\newcommand{\misurau}[3]{\qty{#1}{#2}}
\newcommand{\studio}[1]{\texttt{doc/study/sscdo2/#1/}}
\newcommand{\file}[1]{\texttt{\detokenize{#1}}}
\newcommand{\vuoto}[1]{\ifx\relax#1\relax ---\else #1\fi}

% ---- a card per control
\newcounter{scheda}
\newcommand{\scheda}[8]{%
  \refstepcounter{scheda}\label{scheda:#1}%
  \begin{tcolorbox}[breakable,colback=white,colframe=black!40,boxrule=0.4pt,
    title={\textsf{\bfseries Scheda \thescheda\quad #2}\hfill #7},coltitle=black,colbacktitle=black!6]
  \begin{tabularx}{\linewidth}{@{}>{\bfseries\small\raggedright\arraybackslash}p{3.9cm}X@{}}
  Ambito & \vuoto{#3}\\
  Valore iniziale & \vuoto{#4}\\
  Conseguenze & \vuoto{#5}\\
  Rispetto all'originale & \vuoto{#6}\\
  Studio e audio & \vuoto{#8}\\
  \end{tabularx}
  \end{tcolorbox}}

% ---- a rehearsal card, with blank fields to fill by hand
\newcounter{prova}
\newcommand{\campo}[1]{\par\noindent\textbf{#1:}\ \hrulefill\par\vspace{2mm}\noindent\hrulefill\par}
\newcommand{\prova}[4]{%
  \refstepcounter{prova}\label{prova:#1}%
  \begin{tcolorbox}[breakable,colback=white,colframe=sdaprovare!60,boxrule=0.4pt,
    title={\textsf{\bfseries Prova \theprova}},coltitle=black,colbacktitle=sdaprovare!8]
  \textbf{Ascoltare:} #2\par
  \textbf{Valore:} #3\par
  \textbf{Riferimento:} #4\par\medskip
  \campo{Esito}\campo{Decisione}
  \end{tcolorbox}}

% ---- a numbered question for Davide
\newcounter{questione}
\newcommand{\questione}[2]{%
  \refstepcounter{questione}\label{q:#1}%
  \par\noindent\domanda\ \textbf{Q\thequestione.}\ #2\par\medskip}

% ---- let long monospace names break lines instead of overflowing
\setlength{\emergencystretch}{2em}


\title{Studio sul Corpo d'Ombra \#2\\\large Il porting SEAM: principi, macchine, prove}
\author{Giuseppe Silvi --- SEAM\\\small con Davide Tedesco e Alice Cortegiani}
\date{Versione del \today}

\begin{document}
\maketitle
\begin{abstract}
% abstract.tex
Questo documento raccoglie il porting SEAM di \emph{Studio sul Corpo d'Ombra \#2}.

\end{abstract}
\tableofcontents

\chapter*{Stato delle schede}\label{cap:stati}
\addcontentsline{toc}{chapter}{Stato delle schede}
Ogni scheda della seconda parte, con il suo stato; dopo ogni prova si aggiorna la scheda e questo elenco si rigenera.
\medskip

\listastati

\part{Principi e processi}
% parte1-principi.tex
\chapter{Principi}
Questo capitolo raccoglie i principi del porting.


\part{Le macchine}
% parte2-intro.tex
\chapter{Il sistema}
Questo capitolo descrive il sistema.

% parte2-lmo.tex -- chapter LMO.
\chapter{LMO}\label{cap:lmo}
LMO (\texttt{liveMultiOscs}) è il generatore del brano: una banda stretta di rumore per ciascuno dei quattro canali.
La catena ha tre blocchi.
Il rumore bianco viene da \texttt{no.multinoise} della libreria standard, con un flusso indipendente per ogni canale.
I suoi otto flussi sono i primi due blocchi di un unico \texttt{no.multinoise(12)}, il cui terzo blocco va al coro (scheda~\ref{scheda:coro-voci}).
Il filtro di banda è un passa-alto e un passa-basso di Butterworth di ordine 24 alla stessa frequenza, la riga scritta da Davide.
Due oscillatori, cioè due bande, possono battere tra loro a una distanza $\Delta$ scelta dall'esecutore.
Nella libreria SEAM le tre funzioni sono \texttt{sdt.lmoband}, \texttt{sdt.lmoosc} e \texttt{sdt.lmo(N, f, d)}.

\begin{center}
\begin{tikzpicture}[node distance=5mm and 9mm]
\node[blocco] (n) {rumore bianco, blocchi 1--2\\\texttt{no.multinoise(3N)}};
\node[blocco,above right=2mm and 9mm of n] (b1) {banda a $f-\Delta/2+k$\\HP24 : LP24};
\node[blocco,below right=2mm and 9mm of n] (b2) {banda a $f+\Delta/2+k$\\HP24 : LP24};
\node[blocco,below right=2mm and 9mm of b1] (s) {somma\\$\times 1/\sqrt{2}$};
\node[blocco,right=of s] (v) {volume\\CC81};
\draw[freccia] (n) |- (b1); \draw[freccia] (n) |- (b2);
\draw[freccia] (b1) -| (s); \draw[freccia] (b2) -| (s); \draw[freccia] (s) -- (v);
\end{tikzpicture}\\[1mm]
{\small Un canale $k$ di LMO; i quattro canali sono uguali, con flussi di rumore diversi.}
\end{center}

\scheda{lmo-frequenza}{Frequenza (cue)}
{guidata dai cue; il glissando del cue 2 dura \qty{120}{\second}; nel plugin sono due controlli, f e glide, il tempo del prossimo spostamento di f: un cue imposta prima glide, poi f}
{cue 1: \misurau{97.44}{\hertz}{log}}
{il cue 2 porta la banda da \misurau{97.44}{\hertz}{log} a \misurau{112.67}{\hertz}{log}; la frequenza è il centro della banda di tutti e quattro i canali, a meno dello scarto di $+i$~Hz del canale $i$}
{l'originale moltiplicava la frequenza del cursore per 1,015, un residuo della vecchia formula dei due oscillatori; il porting lo toglie e scrive nei cue le frequenze che Davide sentiva (96 e 111 Hz nell'originale diventano i valori qui sopra); \S\ref{sec:principio-ragiona}}
{\deciso}
{\studio{lmo-bandfilter}, lista dei cue nel log di sessione}

\scheda{lmo-delta}{Distanza $\Delta$ tra i due oscillatori}
{da 0 a \qty{50}{\hertz}, allargato da Giuseppe per sperimentare}
{\qty{0}{\hertz}: le due bande hanno la stessa frequenza e sono decorrelate}
{a distanza $\Delta$ ogni banda si sposta di $\Delta/2$ dal centro; una banda di rumore oscilla già da sola, e il battimento emerge sopra quell'oscillazione da circa 10--15 Hz: l'inviluppo sale di \misura{+3.0}{lmo-beats} dB a 10 Hz, \misura{15.8}{lmo-beats} dB a 20 Hz, \misura{34.5}{lmo-beats} dB a 40 Hz; sotto i 10 Hz la seconda banda allarga il suono senza battere}
{l'originale aveva un secondo oscillatore identico al primo, che raddoppiava il livello; il porting realizza l'intenzione di Davide, due bande che possono battere, con il livello di un oscillatore solo; \S\ref{sec:principio-livelli}}
{\daprovare}
{\studio{lmo-beats}: ascolti a 0, 3, 7, 10, 20, 40 Hz e uno sweep da 0 a 40 Hz}

\scheda{lmo-volume}{Volume (CC81)}
{fader lineare}
{come nel patch}
{a ogni frequenza di campionamento il livello di ogni banda è quello di \qty{96}{\kilo\hertz}: a \qty{48}{\kilo\hertz} il rumore bianco darebbe \misura{3.01}{lmo-beats} dB in più, e LMO lo compensa (\studio{lmo-plugin})}
{nessun guadagno di compatibilità: il livello si fissa in sala; \S\ref{sec:principio-livelli}}
{\daprovare}
{}

\scheda{lmo-canali}{I quattro canali}
{fissi: un flusso di rumore indipendente per canale, scarto di $+i$~Hz sul canale $i$}
{canali decorrelati}
{due canali adiacenti hanno correlazione \misura{-0.007}{lmo-streams}, contro \misura{0.998}{lmo-streams} con un solo flusso condiviso: è questa decorrelazione ad aprire il suono nello spazio; lo scarto di 1 Hz da solo non separa nulla, perché la banda a 1 kHz è larga \misurau{73}{\hertz}{lmo-streams}}
{l'originale ha già un flusso per canale; dal 2 ottobre 2026 gli otto flussi sono i blocchi 1 e 2 di un \texttt{no.multinoise(12)}, di cui il coro usa il blocco 3 (scheda~\ref{scheda:coro-voci}): la statistica resta la stessa, cambia la realizzazione}
{\deciso}
{\studio{lmo-streams}: A/B tra un flusso e più flussi, a quattro canali e in cuffia}

\scheda{lmo-filtro}{Il filtro di banda}
{fisso: passa-alto e passa-basso di ordine 24 alla stessa frequenza}
{versione B}
{la banda a $-3$~dB è larga il \misura{7.35}{lmo-bandfilter}\,\% del centro (Q circa \misura{13.6}{lmo-bandfilter}); la versione C, un passa-banda di Butterworth con gli stessi bordi, ha fianchi quasi verticali ma una coda di \misurau{4.95}{\second}{lmo-bandfilter} contro \misurau{0.31}{\second}{lmo-bandfilter}}
{B è la riga di Davide sulle librerie Faust aggiornate, scelta da Giuseppe all'ascolto tra A (l'originale), B, C e due ordini intermedi, C3 e C2; tutte restano negli studi per il confronto con Davide}
{\daprovare}
{\studio{lmo-bandfilter}: le cinque versioni, con ascolto alla cieca}

% parte2-delrm.tex -- chapter delRM.
\chapter{delRM}\label{cap:delrm}
delRM (\texttt{delRM\_duet}) lavora su quattro canali, e ciascuno elabora soltanto il proprio ingresso dalla TETRAREC.
I canali 1 e 3 sono un filtro a pettine: il segnale più se stesso ritardato di D, $x + x[n-D]$.
I canali 2 e 4 moltiplicano tre segnali, il ritardato, il diretto e il suo integrale, $10 \cdot x[n-D] \cdot x \cdot \int x$, e passano il prodotto a un compressore.
Il ritardo D è quello del plugin DDELAY: una distanza in metri, convertita alla velocità del suono di \misurau{331.4}{\metre\per\second}{delrm-delay} e portata al primo successivo, secondo l'accordo con Davide (\S\ref{sec:principio-genealogia}).
La specifica Faust di DDELAY è verificata valore per valore contro il plugin (\studio{delrm-delay}).
Gli studi usano come segnale di prova una registrazione del Do grave di un clarinetto contrabbasso, con fondamentale a \qty{29.7}{\hertz} e parziali soprattutto dispari, fornita da Giuseppe (appendice~\ref{app:audio}).
Nella libreria SEAM i blocchi sono \texttt{sdt.delrmcomb}, \texttt{sdt.delrmint}, \texttt{sdt.delrmrm} e \texttt{sdt.delrmdyn}.

\begin{center}
\begin{tikzpicture}[node distance=5mm and 8mm]
\node[blocco] (x) {ingresso $x$\\TETRAREC};
\node[blocco,above right=3mm and 8mm of x] (c) {canali 1, 3\\$x + x[n-D]$};
\node[blocco,below right=3mm and 8mm of x] (r) {canali 2, 4\\$10 \cdot x[n-D] \cdot x \cdot \int x$};
\node[blocco,right=of r] (k) {compressore\\11:1};
\node[blocco,right=8mm of k.east,anchor=west,yshift=12mm] (v) {volume\\CC82};
\draw[freccia] (x) |- (c); \draw[freccia] (x) |- (r); \draw[freccia] (r) -- (k);
\draw[freccia] (c) -| (v); \draw[freccia] (k) -| (v);
\end{tikzpicture}\\[1mm]
{\small D è il ritardo primo di DDELAY, regolato in metri.}
\end{center}

\scheda{delrm-distanza}{Distanza (ritardo D)}
{da 0 a \qty{30}{\metre}; si regola durante il montaggio e resta ferma nel brano, e il piccolo click di un cambio è accettato}
{\misurau{7.291}{\metre}{delrm-comb}, cioè \misura{2113}{delrm-comb} campioni a \qty{96}{\kilo\hertz}: i \qty{22}{\milli\second} di Davide}
{il pettine alza e scava i parziali del suono secondo la distanza: a 7,291 m rinforza il terzo parziale del clarinetto contrabbasso di \misura{+6.0}{delrm-comb} dB e scava il settimo di \misura{-7.3}{delrm-comb} dB; a 10 m alza la fondamentale di \misura{+5.5}{delrm-comb} dB e toglie il quinto parziale, \misura{-17.3}{delrm-comb} dB; regolare il ritardo è regolare un timbro}
{l'originale usava secondi con \texttt{de.delay}; il porting usa il ritardo primo di DDELAY, in metri, così lo stesso strumento serve a SSCDO\#4; \S\ref{sec:principio-primi}}
{\daprovare}
{\studio{delrm-comb}: la nota del clarinetto a 5, 7,291 e 10 m, e la nota senza effetto}

\scheda{delrm-volume}{Volume (CC82)}
{fader lineare; secondo la lettura del patch agisce su uno stadio di guadagno di Pure Data}
{da verificare: delRM ha anche un proprio volume interno, che nel wrapper di Pd parte da 0 e nel file \texttt{.dsp} da 0,9}
{nessuna oltre al livello}
{invariato; resta da capire quale dei due volumi il fader muove e con quale valore interno si suonava (domanda~\ref{q:volume-delrm})}
{\domanda}
{log di sessione, ``Build flags''}

\scheda{delrm-dinamica}{Canali 2 e 4: la finestra dinamica}
{dipende dal guadagno dell'ASP880, cioè dal livello con cui il suono entra}
{da trovare in sala}
{il prodotto di tre segnali è cubico: 6 dB in più all'ingresso diventano 18 dB in più prima del compressore; sotto la soglia l'effetto cresce in fretta, sopra il compressore 11:1 lavora come un limitatore; circa 15 dB di dinamica del suonare separano l'effetto assente da quello saturo; i transitori superano l'attacco di 30 ms e toccano 0 dBFS con \misurau{+6}{\decibel}{delrm-rm} all'ingresso, prima del master}
{il compressore è quello dell'originale, con gli stessi parametri}
{\daprovare}
{\studio{delrm-rm}: ascolti del prodotto triplo dopo il compressore}

\scheda{delrm-integratore}{L'integratore}
{fisso}
{---}
{l'integratore originale sommava tutto dall'accensione e trascinava la componente continua: una nota che arrivava dopo cinque minuti produceva un prodotto triplo più forte di \misurau{7.3}{\decibel}{delrm-rm} e simile per \misura{0.88}{delrm-rm} di correlazione a una semplice automodulazione; con l'integratore SEAM l'effetto resta quello dell'accensione per tutto il brano}
{integratore con perdita a 1 Hz, molto al di sotto della fondamentale del clarinetto, scalato per 96000 perché il suono sia quello che Davide sentiva a \qty{96}{\kilo\hertz}; \S\ref{sec:principio-96k}}
{\deciso}
{\studio{delrm-rm}: la nota all'accensione e dopo cinque minuti, nell'originale e nel porting}

\scheda{delrm-dcblocker}{I DC blocker}
{due in serie nell'originale: uno alla fine di ogni canale, uno dopo il volume}
{assenti nel porting}
{con l'integratore SEAM non c'è componente continua da trattenere; a \qty{96}{\kilo\hertz} i DC blocker erano passa-alto a \misurau{76.59}{\hertz}{delrm-dcblock}, e in serie toglievano \misurau{17.63}{\decibel}{delrm-dcblock} alla fondamentale del clarinetto: il basso della performance era più sottile}
{tolti perché il loro problema non esiste più (\S\ref{sec:principio-processo}); se Davide vuole quel basso più sottile, rientra come un passa-alto dichiarato, due \texttt{fi.dcblockerat(76.59)} in serie, come erano i due DC blocker}
{\domanda}
{\studio{delrm-dcblock}: A/B della nota senza e con i due DC blocker}

% parte2-stunedrev.tex -- chapter stunedrev.
\chapter{stunedrev}\label{cap:stunedrev}
stunedrev, l'APF della partitura, è fatto di quattro linee indipendenti, una per ogni faccia di STONED.
Ogni linea è una catena di 42 filtri passa-tutto (all-pass) in serie, nella forma di Moorer, con guadagno $g = 1/\sqrt{2}$ (\studio{stunedrev-allpass}).
La sezione $i$ ritarda il suono di un tempo pari a quello del cursore moltiplicato per $(i+1) \cdot k$ e portato al primo successivo, con $k = \sqrt{2}$, $\varphi$, $e$, $\pi$ per le quattro facce (\studio{stunedrev-delays}).
La macchina viene da in-phi-rev di \emph{Canto alla durata} (\S\ref{sec:principio-genealogia}).
Nella libreria SEAM le funzioni sono \texttt{sdt.stdel}, \texttt{sdt.stmd}, \texttt{sdt.stline} e \texttt{sdt.stunedrev}.

\section{Una memoria lunga}
Un filtro all-pass lascia passare tutta l'energia che riceve, ma la restituisce nel tempo.
Ogni sezione ritarda l'energia in media esattamente del proprio tempo $t$, qualunque sia $g$.
In una catena le medie si sommano: l'energia di una linea arriva in media dopo la somma dei suoi 42 ritardi.
Il percorso diretto attraversa le 42 sezioni con il guadagno $-g$ ciascuna e si riduce a $-126$~dB: nei primi istanti dopo l'ingresso di un suono l'uscita è quasi muta, e il suono riemerge man mano che i ritardi si sommano.
stunedrev è dunque una memoria più che un riverbero: ciò che entra torna, diffuso, dopo decine di secondi o minuti.
Giuseppe e Davide lo hanno voluto così: ogni faccia di STONED restituisce il suono con la propria scala del tempo.

\scheda{stunedrev-tempi}{I quattro tempi ($\sqrt{2}$, $\varphi$, $e$, $\pi$)}
{da 1 a \qty{100}{\milli\second} per ogni linea}
{83, 47, 7 e \qty{71}{\milli\second}, come nel patch Pure Data (il valore 33 del file \texttt{.dsp} è solo il default del cursore)}
{il baricentro dell'energia cade a \misurau{106.0}{\second}{stunedrev-chain} per $\sqrt{2}$, \misurau{68.7}{\second}{stunedrev-chain} per $\varphi$, \misurau{17.2}{\second}{stunedrev-chain} per $e$ e \misurau{201.4}{\second}{stunedrev-chain} per $\pi$; una nota di clarinetto torna al 99,9\,\% della sua energia dopo \misurau{327.8}{\second}{stunedrev-chain}, \misurau{171.1}{\second}{stunedrev-chain}, \misurau{44.4}{\second}{stunedrev-chain} e \misurau{603.8}{\second}{stunedrev-chain}; il baricentro cresce in proporzione al tempo del cursore, perché è il tempo moltiplicato per $k$ e per 903, la somma dei fattori da 1 a 42}
{la struttura è quella di Davide; i primi seguono la regola SEAM (scheda~\ref{scheda:stunedrev-sezioni})}
{\daprovare}
{\studio{stunedrev-chain}: la nota del clarinetto nelle quattro linee, da rigenerare con \texttt{render.py}}

\scheda{stunedrev-ingresso-uscita}{Ingresso (CC83) e uscita (CC84)}
{fader lineari}
{come nel patch}
{nel patch Pure Data un \texttt{adc\textasciitilde{} 5 6 7 8} è collegato anche direttamente all'ingresso del riverbero; Pd somma i segnali su un ingresso, quindi stunedrev riceve sempre il suono e il fader APF INPUT aggiunge al massimo 6~dB}
{il collegamento diretto sembra un resto delle prove, ma riprodurlo o toglierlo è una scelta di Davide}
{\domanda}
{log di sessione, ``Double feed into the reverb''}

\scheda{stunedrev-sezioni}{Le sezioni che cambiano}
{fisse: dipendono dai quattro tempi}
{7 sezioni su 168 diverse dall'originale ai tempi iniziali}
{il porting arrotonda i campioni dove Davide li troncava: sui 16\,800 ritardi possibili, \misura{813}{stunedrev-delays} passano a un altro primo, e lo spostamento arriva a \misura{96}{stunedrev-delays} campioni, un millisecondo; ai tempi iniziali cambiano 3 sezioni di $\sqrt{2}$, nessuna di $\varphi$, 3 di $e$ e 1 di $\pi$, ciascuna di un solo primo; dove i primi coincidono la linea è identica a quella di Davide campione per campione}
{regola SEAM mantenuta per coerenza con il sistema; \S\ref{sec:principio-primi}}
{\deciso}
{\studio{stunedrev-delays}: la tabella delle sezioni, e il confronto con l'originale}

\scheda{stunedrev-memoria}{Memoria}
{fissa}
{---}
{l'originale in Faust occupa \misura{3.94}{stunedrev-chain} GiB, l'external di Pure Data \misura{15.2}{stunedrev-chain} GiB (da qui la nota di Davide sui 32 GB di RAM); la versione SEAM in Faust occupa \misura{1.66}{stunedrev-chain} GiB, oppure \misura{1.15}{stunedrev-chain} GiB con l'opzione \texttt{-dlt 4096}; il plugin C++ dimensionato a \qty{96}{\kilo\hertz} occuperà \misura{588}{stunedrev-chain} MiB}
{nell'originale tutte le sezioni di una linea avevano il buffer della più lunga; nel porting ogni sezione ha il buffer del proprio ritardo massimo, e il suono resta identico}
{\deciso}
{\studio{stunedrev-chain}}

% parte2-coro-master.tex -- choir placeholder and master.
\chapter{Coro e master}\label{cap:coro-master}
Il coro (\texttt{pitchDetectorChoirMcAdams}) suona in quattro istanze, una per canale.
Ogni istanza ha 16 bande: l'inviluppo della banda $k$ intorno a $f \cdot k$ modula un rumore filtrato intorno a $f \cdot k^{a}$, con $f$ = 48, 48, 96 e \qty{96}{\hertz}, $a$ = 1; 1,01; 1,1; 0,9, fattore di merito 350 e rilascio di \qty{1.5}{\second}.
Lo studio del coro seguirà quello di stunedrev: la scheda qui sotto tiene il suo posto nel sistema.

\scheda{coro-uscita}{Uscita del coro (CC86)}
{}
{}
{non ancora studiato: lo studio seguirà quello di stunedrev}
{}
{\daprovare}
{}

\section{Master}
Il master chiude la catena e porta i quattro canali alle uscite 9--12, verso STONED.

\scheda{master}{Master (CC88)}
{fader lineare, levigato da \texttt{si.smoo} della libreria standard}
{come nel patch}
{nessuna oltre al livello}
{invariato; nel porting gli stadi di guadagno del patch diventano i fader delle tracce di Reaper}
{\daprovare}
{log di sessione, blocco 6 di delRM}


\part{Protocollo di prova}
% parte3-prove.tex
\chapter{Schede di prova}
Questo capitolo raccoglie le schede di prova.


\appendix
% appendice-audio.tex -- index of the audio files.
\chapter{File audio}\label{app:audio}
Ogni studio con un suono ha una cartella \texttt{renders/} con i file audio e un \file{render-log.md} che dice come sono stati prodotti: livello, durata, frequenza di campionamento.
I file brevi sono nel repository; quelli di stunedrev, lunghi, si rigenerano con \file{render.py} (capitolo~\ref{cap:preparazione}).
La sorgente delle prove su strumento è il Do grave di un clarinetto contrabbasso, registrato a \qty{96}{\kilo\hertz} con il microfono a un metro, pubblicato con l'accordo di Giuseppe.

\section*{LMO: il filtro di banda}
\studio{lmo-bandfilter}, nel repository.
\begin{longtable}{@{}>{\raggedright\arraybackslash}p{6.2cm}>{\raggedright\arraybackslash}p{9.2cm}@{}}
\toprule
File & Contenuto\\
\midrule
\file{steady96_A.wav} \ldots\ \file{steady96_C3.wav} & la banda ferma a \qty{97.44}{\hertz} nelle versioni A, B, C, C2, C3\\
\file{steady96_A_davide_level.wav} & la versione A al livello reale dell'originale, raddoppiato dai due oscillatori identici\\
\file{gliss_A.wav} \ldots\ \file{gliss_C3.wav} & il glissando del cue 2 nelle cinque versioni\\
\file{step_A.wav} \ldots\ \file{step_C3.wav} & un salto di frequenza, per sentire il transitorio\\
\file{blind/*_V.wav} \ldots\ \file{blind/*_Z.wav} & gli stessi tre ascolti con le versioni nascoste; la chiave è in \file{blind/KEY.txt}\\
\bottomrule
\end{longtable}

\section*{LMO: due oscillatori che battono}
\studio{lmo-beats}, nel repository.
\begin{longtable}{@{}>{\raggedright\arraybackslash}p{6.2cm}>{\raggedright\arraybackslash}p{9.2cm}@{}}
\toprule
File & Contenuto\\
\midrule
\file{one_osc_mono.wav} & un oscillatore solo, per confronto\\
\file{d00_mono.wav} \ldots\ \file{d40_mono.wav} & due bande a $\Delta$ = 0, 3, 7, 10, 20, \qty{40}{\hertz}, in mono\\
\file{d00_4ch.wav} \ldots\ \file{d40_4ch.wav} & le stesse a quattro canali\\
\file{sweep_d00-40_mono.wav} & $\Delta$ che sale da 0 a \qty{40}{\hertz}\\
\bottomrule
\end{longtable}

\section*{LMO: uno o più flussi di rumore}
\studio{lmo-streams}, nel repository.
\begin{longtable}{@{}>{\raggedright\arraybackslash}p{6.2cm}>{\raggedright\arraybackslash}p{9.2cm}@{}}
\toprule
File & Contenuto\\
\midrule
\file{lmo_gs1_4ch.wav}, \file{lmo_gs1_ch01_stereo.wav} & la variante di Giuseppe con un flusso, a quattro canali e in cuffia\\
\file{lmo_gs2_4ch.wav}, \file{lmo_gs2_ch01_stereo.wav} & la stessa con due flussi\\
\file{lmo_dav_4ch.wav}, \file{lmo_dav_ch01_stereo.wav} & un flusso per canale, come nell'originale e nel porting\\
\bottomrule
\end{longtable}

\section*{delRM: il pettine}
\studio{delrm-comb}, nel repository.
\begin{longtable}{@{}>{\raggedright\arraybackslash}p{6.2cm}>{\raggedright\arraybackslash}p{9.2cm}@{}}
\toprule
File & Contenuto\\
\midrule
\file{ccb_dry.wav} & la nota del clarinetto senza effetto; è anche la sorgente dei render di stunedrev\\
\file{ccb_comb_5.000m.wav}, \file{ccb_comb_7.291m.wav}, \file{ccb_comb_10.000m.wav} & il pettine a tre distanze\\
\bottomrule
\end{longtable}

\section*{delRM: il prodotto triplo}
\studio{delrm-rm}, nel repository.
\begin{longtable}{@{}>{\raggedright\arraybackslash}p{6.2cm}>{\raggedright\arraybackslash}p{9.2cm}@{}}
\toprule
File & Contenuto\\
\midrule
\file{note_dry.wav} & la nota senza effetto\\
\file{post_I0-original_fresh.wav}, \file{post_I0-original_after5min.wav} & l'integratore originale, all'accensione e dopo cinque minuti\\
\file{post_I1-leaky-1Hz-96k_fresh.wav}, \file{post_I1-leaky-1Hz-96k_after5min.wav} & l'integratore SEAM, all'accensione e dopo cinque minuti\\
\file{post_I2-dcblocked_fresh.wav}, \file{post_I2-dcblocked_after5min.wav} & un'alternativa scartata, con un DC blocker prima dell'integratore\\
\file{selfrm_post.wav} & la semplice automodulazione, per confronto\\
\bottomrule
\end{longtable}

\section*{delRM: i DC blocker}
\studio{delrm-dcblock}, nel repository.
\begin{longtable}{@{}>{\raggedright\arraybackslash}p{6.2cm}>{\raggedright\arraybackslash}p{9.2cm}@{}}
\toprule
File & Contenuto\\
\midrule
\file{note_dry.wav} & la nota senza effetto\\
\file{ch1_comb_clean.wav}, \file{ch1_comb_2dcblockers.wav} & il pettine senza e con i due DC blocker\\
\file{ch2_triple_clean.wav}, \file{ch2_triple_2dcblockers.wav} & il prodotto triplo senza e con i due DC blocker\\
\bottomrule
\end{longtable}

\section*{stunedrev: le quattro facce}
\studio{stunedrev-chain}, da rigenerare con \file{render.py}; nel repository c'è solo \file{render-log.md}.
\begin{longtable}{@{}>{\raggedright\arraybackslash}p{6.2cm}>{\raggedright\arraybackslash}p{9.2cm}@{}}
\toprule
File & Contenuto\\
\midrule
\file{note_dry.wav} & la nota senza effetto\\
\file{stunedrev_sqrt2_83ms.wav} & la faccia $\sqrt{2}$ a \qty{83}{\milli\second}, \qty{327.8}{\second}\\
\file{stunedrev_phi_47ms.wav} & la faccia $\varphi$ a \qty{47}{\milli\second}, \qty{171.1}{\second}\\
\file{stunedrev_e_7ms.wav} & la faccia $e$ a \qty{7}{\milli\second}, \qty{44.4}{\second}\\
\file{stunedrev_pi_71ms.wav} & la faccia $\pi$ a \qty{71}{\milli\second}, \qty{603.8}{\second}\\
\bottomrule
\end{longtable}

\end{document}
